\section{The Lean companion and its limits}
The file \code{proofs.lean} contains the complete small examples and exercise solutions from Chapters~15--16, within the namespace \code{MathCourse}. The file \code{lean-toolchain} pins Lean 4.19.0. It does not pin the toolchain of the external research repository; that project must use its own configuration and dependencies.

After installing the official Lean toolchain support, run
\par\noindent\begin{minipage}{\linewidth}
\begin{lstlisting}
lean proofs.lean
python verify.py --lean
\end{lstlisting}
\end{minipage}\par

The second command requires Lean execution in addition to the numerical tests. It fails rather than silently reporting a pass when Lean is unavailable. Inspect the produced \code{lean\_check.log} and the theorem dependency output. The intentionally impossible-context theorem remains a conditional statement, not an unconditional contradiction, even if it uses no extra axiom.

\textbf{Execution status of this edition.} The Lean examples have been reviewed against their mathematical arguments and the documented language constructs, but no Lean executable was available in the authoring environment, and a toolchain could not be obtained there. They were therefore \emph{not compiler-checked in this edition}. The companion is supplied for execution with the pinned toolchain; its acceptance is not represented as an already completed test. The full OpenAI research repository was also not built or independently verified for this book.

For additional practice after these lessons, \emph{Mathematics in Lean} provides an interactive textbook and exercises~\cite{mil}. Use it as a continuation once the statement and proof-state ideas in this course are familiar, not as an unexplained prerequisite imposed before the first lesson.
